% ===== PRÉAMBULE CANONIQUE — à coller VERBATIM en tête de CHAQUE .tex =====
\documentclass[11pt,a4paper]{article}
\usepackage[utf8]{inputenc}\usepackage[T1]{fontenc}\usepackage{lmodern}
\usepackage[french]{babel}
\usepackage{amsmath,amssymb,mathtools}\usepackage{icomma}
\usepackage[version=4]{mhchem}          % \ce{...} : équations & formules
\usepackage{siunitx}\sisetup{locale=FR} % \qty{}{}, \si{}, \ang{}
\usepackage{chemfig}                    % \chemfig{...} : structures développées
\usepackage{xcolor}\usepackage{colortbl}
\usepackage{tikz}\usepackage{pgfplots}\pgfplotsset{compat=1.18}
\usepackage[most]{tcolorbox}\usepackage{enumitem}\usepackage{array,booktabs}
\usepackage[a4paper,margin=2.2cm]{geometry}
\usetikzlibrary{arrows.meta,calc,angles,quotes,positioning,patterns,backgrounds,shapes.geometric}
\definecolor{primary}{RGB}{222,49,99}\definecolor{primarydark}{RGB}{150,22,55}
\definecolor{defcol}{RGB}{0,114,178}\definecolor{methcol}{RGB}{52,92,148}
\definecolor{excol}{RGB}{0,135,90}\definecolor{attncol}{RGB}{200,40,55}
\tcbset{boxbase/.style={enhanced,breakable,boxrule=0.9pt,arc=2.5pt,
  left=3mm,right=3mm,top=2.3mm,bottom=2.3mm,fonttitle=\bfseries,coltitle=white}}
% --- boîtes TUTORIEL (noms EXACTS, ne pas renommer) ---
\newtcolorbox{cle}[1][]{boxbase,colback=defcol!6,colframe=defcol,
  title={\textsc{À retenir}\ifx\relax#1\relax\relax\else\textmd{ : \itshape #1}\fi}}
\newtcolorbox{methode}[1][]{boxbase,colback=methcol!7,colframe=methcol,sharp corners,
  title={\textsc{Méthode}\ifx\relax#1\relax\relax\else\textmd{ --- \itshape #1}\fi}}
\newtcolorbox{exemplebox}[1][]{boxbase,colback=excol!5,colframe=excol,
  title={\textsc{Exemple}\ifx\relax#1\relax\relax\else\textmd{ : \itshape #1}\fi}}
\newtcolorbox{piege}[1][]{boxbase,colback=attncol!6,colframe=attncol,
  title={\textsc{Piège}\ifx\relax#1\relax\relax\else\textmd{ : \itshape #1}\fi}}
\newtcolorbox{remarque}[1][]{boxbase,colback=black!4,colframe=black!45,
  title={\textsc{Remarque}\ifx\relax#1\relax\relax\else\textmd{ : \itshape #1}\fi}}
% --- boîtes EXERCICE / CORRIGÉ (noms EXACTS, auto-comptées) ---
\newtcolorbox[auto counter]{exo}[1][]{boxbase,colback=primary!4,colframe=primary,
  title={\textsc{Exercice}~\thetcbcounter\ifx\relax#1\relax\relax\else\textmd{ --- \itshape #1}\fi}}
\newtcolorbox[auto counter]{corrige}[1][]{boxbase,colback=excol!4,colframe=excol,
  title={\textsc{Corrigé}~\thetcbcounter\ifx\relax#1\relax\relax\else\textmd{ --- \itshape #1}\fi}}
\newcommand{\R}{\mathbb{R}}\newcommand{\N}{\mathbb{N}}\newcommand{\Z}{\mathbb{Z}}\newcommand{\ds}{\displaystyle}
% (un \usepackage{fancyhdr}+en-tête est possible ; il est ignoré côté web)
% =========================================================================
\begin{document}
\newcommand{\CF}{\ensuremath{\mathrm{C.F.}}}
\begin{corrige}[deux structures de \ce{ClO4-}]
\begin{enumerate}[label=\alph*)]
  \item (A) \ce{Cl} : $7 - 0 - 4 = \mathbf{+3}$ ; chaque \ce{O} : $6 - 6 - 1 = \mathbf{-1}$.
        Somme : $+3 + 4\times(-1) = -1$. \checkmark
  \item (B) \ce{Cl} : $7 - 0 - 7 = \mathbf{0}$ ($3$ doubles $+$ $1$ simple $= 7$ liaisons) ;
        \ce{O} double : $6 - 4 - 2 = \mathbf{0}$ ; \ce{O} simple : $6 - 6 - 1 = \mathbf{-1}$.
        Somme : $0 + 3\times 0 + (-1) = -1$. \checkmark
  \item La structure \textbf{(B)} est la plus représentative : ses charges formelles sont
        minimales ($\CF(\ce{Cl})=0$ au lieu de $+3$). Le chlore, en \textbf{$3^{\text{e}}$
        période}, a le droit de dépasser l'octet ($14$~e$^-$), ce qui rend (B) permise.
\end{enumerate}
\end{corrige}

\begin{corrige}[la meilleure structure de \ce{SOCl2}]
\begin{enumerate}[label=\alph*)]
  \item \ce{S=O} double — \ce{S} ($1$ doublet, $4$ liaisons : $2+1+1$) : $6-2-4 = \mathbf{0}$ ;
        \ce{O} : $6-4-2 = \mathbf{0}$ ; chaque \ce{Cl} : $7-6-1 = \mathbf{0}$.
  \item \ce{S-O} simple — \ce{S} ($1$ doublet, $3$ liaisons) : $6-2-3 = \mathbf{+1}$ ;
        \ce{O} : $6-6-1 = \mathbf{-1}$ ; chaque \ce{Cl} : $\mathbf{0}$.
  \item On retient la structure à \textbf{\ce{S=O} double} : toutes ses charges formelles sont
        nulles, alors que la structure simple fait apparaître des charges $\pm 1$ inutiles. Le
        soufre ($3^{\text{e}}$ période) peut porter cet octet étendu.
\end{enumerate}
\end{corrige}

\begin{corrige}[toutes les charges de \ce{HNO3}]
\begin{enumerate}[label=\alph*)]
  \item \ce{N} : $5 - 0 - 4 = \mathbf{+1}$.
  \item \ce{O} doublement lié : $6 - 4 - 2 = \mathbf{0}$.
  \item \ce{O} du groupe \ce{-O-H} : $6 - 4 - 2 = \mathbf{0}$.
  \item \ce{O} simplement lié seul : $6 - 6 - 1 = \mathbf{-1}$.
  \item Somme : $(+1) + 0 + 0 + (-1) + \underbrace{0}_{\ce{H}} = \mathbf{0}$ : la molécule est
        bien neutre. \checkmark
\end{enumerate}
\end{corrige}

\begin{corrige}[quand les charges formelles ne suffisent pas]
\begin{enumerate}[label=\alph*)]
  \item \textbf{Non} : A et B ont les mêmes charges formelles en valeur absolue ($+1$ et $-1$).
        Le critère « \CF{} minimales » les met à égalité.
  \item On applique le critère de l'\textbf{électronégativité} : la charge \textbf{négative} doit
        se placer sur l'atome le plus électronégatif. Comme \ce{O} est plus électronégatif que
        \ce{C}, on retient la structure \textbf{A} (le $-1$ sur l'oxygène).
\end{enumerate}
\end{corrige}

\begin{corrige}[l'octet étendu améliore-t-il \ce{SO4^2-} ?]
\begin{enumerate}[label=\alph*)]
  \item $4$ liaisons simples — \ce{S} : $6 - 0 - 4 = \mathbf{+2}$ ; chaque \ce{O} :
        $6 - 6 - 1 = \mathbf{-1}$. La \CF{} du soufre est élevée ($+2$).
  \item $2$ liaisons \ce{S=O} — \ce{S} ($0$ doublet, $6$ liaisons) : $6 - 0 - 6 = \mathbf{0}$ ;
        les $2$ \ce{O} doublement liés : $6 - 4 - 2 = \mathbf{0}$ ; les $2$ \ce{O} simples restent
        à $-1$.
  \item Oui, elle est meilleure : $\CF(\ce{S})$ passe de $+2$ à $0$ (charges minimales), possible
        car \ce{S} est en $3^{\text{e}}$ période. Somme : $0 + 2\times 0 + 2\times(-1) = \mathbf{-2}$,
        toujours la charge de l'ion. \checkmark
\end{enumerate}
\end{corrige}

\end{document}
